
\section{Appendix D. Negative Results}

\subsection{D.1 Attempts That Failed}

\subsubsection{D.1.1 Very Small Sketch (r=4)}

\textbf{Goal.} Maximize FLOPs reduction.

\textbf{Result.} Score correlation dropped to {{CORR4}}. Accuracy collapsed.

\textbf{Conclusion.} r=8 is the practical minimum.

\subsubsection{D.1.2 Shared Projection Across Heads}

\textbf{Goal.} Use the same $\Omega$ for all heads.

\textbf{Result.} 0.5\% accuracy drop, no speed benefit.

\textbf{Conclusion.} Head-wise is better.

\subsubsection{D.1.3 Learned Projection}

\textbf{Goal.} Train $\Omega$ jointly with LoRA.

\textbf{Result.} Trained $\Omega$ lost the JL property, accuracy dropped.

\textbf{Conclusion.} $\Omega$ must be random and fixed.

\subsubsection{D.1.4 Sparse Random Projection (Achlioptas)}

\textbf{Goal.} Use Achlioptas-style sparse projections.

\textbf{Result.} Similar accuracy, 10\% faster sketch.

\textbf{Conclusion.} Marginal gain, complexity not worth it.

\subsubsection{D.1.5 Larger Scale Factor}

\textbf{Goal.} Use $1/\sqrt{d_h}$ instead of $1/\sqrt{r}$.

\textbf{Result.} Softmax saturated, accuracy collapsed.

\textbf{Conclusion.} $1/\sqrt{r}$ is correct.

\subsection{D.2 Abandoned Variants}

\subsubsection{D.2.1 Sketching V}

\textbf{Goal.} Also sketch $V$.

\textbf{Result.} Output shape mismatch; problematic for downstream.

\subsubsection{D.2.2 Per-Token Projection}

\textbf{Goal.} Different $\Omega$ per token.

\textbf{Result.} Not possible without breaking JL.

\subsubsection{D.2.3 Hierarchical Sketch}

\textbf{Goal.} Multi-level sketches.

\textbf{Result.} Complexity not worth it.

\subsection{D.3 Failed Hyperparameters}

\subsubsection{Sketch Rank}

\begin{table}[H]
\centering
\small
\begin{tabular}{lll}
\toprule
\textbf{$r$} & \textbf{DistilBERT} & \textbf{Correlation} \\
\midrule
4 & {{D\_R4}} & {{CORR4}} \\
8 & {{D\_R8}} & {{CORR8}} \\
16 & {{D\_R16}} & {{CORR16}} \\
32 & {{D\_R32}} & {{CORR32}} \\
64 & {{D\_R64}} & 1.00 \\
\bottomrule
\end{tabular}
\end{table}

\subsubsection{LoRA Rank}

\begin{table}[H]
\centering
\small
\begin{tabular}{ll}
\toprule
\textbf{$r_L$} & \textbf{DistilBERT} \\
\midrule
0 & {{D\_LORA0}} \\
8 & {{D\_LORA8}} \\
16 & {{D\_LORA16}} \\
32 & {{D\_LORA32}} \\
\bottomrule
\end{tabular}
\end{table}

\subsection{D.4 Multi-Seed}

\begin{table}[H]
\centering
\small
\begin{tabular}{ll}
\toprule
\textbf{Seed} & \textbf{DistilBERT} \\
\midrule
42 & {{SEED42}} \\
123 & {{SEED123}} \\
456 & {{SEED456}} \\
Mean & {{MEAN}} \\
Std & {{STD}} \\
\bottomrule
\end{tabular}
\end{table}

\textbf{Observation.} Seed variance is small (0.1\%) because JL is 
concentration-of-measure.

\subsection{D.5 Lessons Learned}

\begin{enumerate}
\item \textbf{r=8 is the practical minimum.}
\item \textbf{Head-wise projections are better.}
\item \textbf{$\Omega$ must be fixed (not learned).}
\item \textbf{Scale factor must be $1/\sqrt{r}$.}
\item \textbf{RPC does not reduce memory.}

\end{enumerate}
\subsection{D.6 Reproducibility Note}

The random seed for \texttt{proj} initialization affects results. We use seed 
{{SEED}} for all experiments.

